\documentclass[11pt]{article}
\usepackage[utf8]{inputenc}
\usepackage[margin=1in]{geometry}
\usepackage{amsmath}
\usepackage{amsfonts}
\usepackage{amssymb}
\usepackage{mathtools}
\usepackage{booktabs}
\usepackage{tabularx}
\usepackage{microtype}
\usepackage{needspace}
\usepackage{hyperref}
\usepackage{cleveref}

\setcounter{secnumdepth}{3}
\setcounter{tocdepth}{2}
\setlength{\parskip}{0.5\baselineskip}
\setlength{\parindent}{0pt}
\setlength{\tabcolsep}{5pt}
\newcolumntype{Y}{>{\raggedright\arraybackslash}X}
\widowpenalty=10000
\clubpenalty=10000
\displaywidowpenalty=10000
\allowdisplaybreaks[2]
\setlength{\abovedisplayskip}{8pt plus 2pt minus 2pt}
\setlength{\belowdisplayskip}{8pt plus 2pt minus 2pt}
\setlength{\abovedisplayshortskip}{6pt plus 2pt minus 2pt}
\setlength{\belowdisplayshortskip}{6pt plus 2pt minus 2pt}

\hypersetup{
    colorlinks=true,
    linkcolor=blue,
    citecolor=blue,
    urlcolor=blue
}

\title{Multibody Dynamics Formalisms:\\Drift, Control, and the Geometry of Physical Law\\[0.5em]
\large A Unified Treatment of Affine Decomposition Across Calculational Frameworks}
\author{}
\date{}

\begin{document}
\maketitle
\tableofcontents
\newpage

%% ============================================================
\section{Introduction}\label{sec:introduction}
%% ============================================================

\subsection{The Central Question}

Every mechanical system evolves according to physical laws that exist independently of how we choose to describe them.
A double pendulum swinging under gravity does not care whether the analyst uses Lagrangian mechanics, Newton--Euler body-by-body balance, screw-theoretic spatial algebra, or coordinate-free differential geometry.
The underlying physics---the forces, the accelerations, the energy exchanges---remain identical.

Yet different calculational frameworks exist because they illuminate different aspects of the same physics.
Euler--Lagrange methods reveal the energetic structure of motion.
Newton--Euler recursions expose internal load paths between bodies.
Screw-theoretic formulations provide a unified language for spatial velocity and force.
Operational-space methods connect joint-level effort to task-level objectives.
Coordinate-free geometric methods strip away chart-dependent artifacts to reveal the intrinsic structure of dynamics.

This article has three interwoven goals:

\begin{enumerate}
\item \textbf{To demonstrate that every framework admits a clean affine decomposition into drift and control}, and to show precisely what these components look like in each representation.

\item \textbf{To establish that constraint forces and torques, despite doing no net generalized work, are active and essential participants in energy transfer and motion redirection}, and that their inclusion refines but never breaks the drift--control decomposition.

\item \textbf{To make explicit the symmetry and representation-independence that underlies all of mechanics}: the physics does not depend on coordinate labels, reference frames, or the mathematical language chosen for description.
\end{enumerate}

The organizing mathematical structure is the \textbf{control-affine decomposition}:
\begin{equation}\label{eq:affine-master}
\dot{x} = f_{\mathrm{drift}}(x) + G(x)\,u.
\end{equation}
Here $x$ is the system state, $f_{\mathrm{drift}}(x)$ is the drift vector field determined entirely by the current state, and $G(x)\,u$ is the control contribution driven by external inputs.
This decomposition is not a modeling convenience---it is a structural property of mechanical systems with independent actuation channels.


\subsection{Theme 1: The Independence of Drift and Control}

The drift field $f_{\mathrm{drift}}(x)$ encodes everything the system would do on its own---the inertial coupling, the velocity-dependent (Coriolis and centripetal) effects, and gravitational loading.
It is a complete dynamical prediction conditioned on the current state with zero actuation.

The control field $G(x)\,u$ encodes the incremental effect of actuation.
It depends on the current state (through the input map $G(x)$) and on the applied input $u$, but it is structurally independent of the drift field.
This independence means:

\begin{itemize}
\item \textbf{Drift can be evaluated without knowledge of control intent.}
At any instant, one can compute what the system would do with zero input.
This is the \textbf{Zero Torque Counterfactual (ZTCF)}---an instantaneous diagnostic of accumulated state effects.

\item \textbf{Control can be evaluated as a perturbation from drift.}
The actual acceleration is always drift plus the incremental control contribution.
No cross-coupling terms between drift and control exist in a properly decomposed model.

\item \textbf{The decomposition persists across all frameworks.}
Whether we write the equations in generalized coordinates, spatial wrenches, or intrinsic geometric form, the same structural split exists.
\end{itemize}


\subsection{Theme 2: Constraint Forces Are Not Passive}

A widespread misconception treats constraint forces as passive because they do no net generalized work: $P_{\mathrm{constr}} = (J_c^T \lambda)^T \dot{q} = \lambda^T J_c \dot{q} = 0$ when the velocity-level constraint $J_c \dot{q} = 0$ is satisfied.

But ``no net work'' does not mean ``no physical effect.''
Constraint forces actively redirect motion and mediate internal energy transfer between bodies.
In a whip crack, the constraint forces at each link junction are enormous and are the primary mechanism by which kinetic energy migrates from proximal to distal segments.

When constraints are present, the drift--control decomposition must be refined to include constraint reactions in the passive baseline:
\begin{equation}\label{eq:constrained-master}
M\ddot{q} + h(q,\dot{q}) = \tau + J_c^T\lambda + \tau_{\mathrm{ext}}.
\end{equation}
The multiplier $\lambda$ depends on the state and on the actuation, so the ``passive'' baseline is the solution at $\tau = 0$:
\begin{equation}
\ddot{q}_{\mathrm{drift},t}^{c} = \left[M^{-1}\bigl(-h + J_c^T\lambda_0 + \tau_{\mathrm{ext}}\bigr)\right]_{(q(t),\dot{q}(t))}
\end{equation}
where $\lambda_0$ is the constraint multiplier under zero actuation.


\subsection{Theme 3: Symmetry and Representation Independence}

The equations of motion are \textbf{objective}: they express physical law that does not depend on the observer's choice of coordinates, reference frame, or mathematical framework.

This is not merely a philosophical statement---it has concrete computational consequences:
\begin{itemize}
\item \textbf{Frame changes are similarity transformations.}
If $g \in SE(3)$ maps between frames, then twists transform as $V_s = \mathrm{Ad}_g\, V_b$ and wrenches as $F_s = \mathrm{Ad}_g^{-T} F_b$.
Instantaneous power $F^T V$ is invariant.

\item \textbf{Every framework computes the same physical accelerations.}
The Euler--Lagrange $\ddot{q}$, the Newton--Euler body accelerations, the screw-theoretic twist rates, and the operational-space task accelerations are all related by known kinematic transformations.

\item \textbf{Conservation laws are symmetry consequences.}
By Noether's theorem: translational symmetry yields linear momentum conservation, rotational symmetry yields angular momentum conservation, and time-translation symmetry yields energy conservation.
\end{itemize}


\subsection{The Zero Torque Counterfactual (ZTCF)}

The ZTCF is an evaluation rule, not a separate physical concept.
At time $t$, with state $x(t)$:
\begin{equation}\label{eq:ztcf-def}
\dot{x}_{\mathrm{ZTCF},t} \coloneqq \dot{x}\bigl(x(t), u = 0\bigr) = f_{\mathrm{drift}}\bigl(x(t)\bigr).
\end{equation}

ZTCF is simply the drift field sampled at the current state.
It answers the question: ``What would the system do right now if all actuation were removed?''

\textbf{ZTCF is time-local.}
It is defined at a single instant, not as a trajectory prediction.
The first-order propagation
\begin{equation}
x(t + \Delta t) = x(t) + \Delta t\, \mathcal{F}(x(t), u(t)) + \mathcal{O}(\Delta t^2)
\end{equation}
shows that ZTCF is the $u = 0$ evaluation of a local vector field, not a statement about global future behavior.

\textbf{ZTCF in constrained systems.}
When constraints are present, ``zero actuation'' still requires constraint reactions for kinematic compatibility:
\begin{equation}
\ddot{q}_{\mathrm{ZTCF},t}^{c} = \ddot{q}_{\mathrm{drift},t}^{c},
\end{equation}
where the constrained drift includes $J_c^T \lambda_0$, the reaction forces under zero actuation.


\subsection{State Accumulation and Causality}

The current state stores the entire past acceleration history:
\begin{equation}\label{eq:state-accumulation}
\dot{q}(t) = \dot{q}(t_0) + \int_{t_0}^{t} \ddot{q}(\sigma)\,d\sigma, \qquad h_m(t) = M(q(t))\,\dot{q}(t),
\end{equation}
and in first-order state form:
\begin{equation}
x(t) = x(t_0) + \int_{t_0}^{t} \mathcal{F}(x(\sigma), u(\sigma))\,d\sigma.
\end{equation}

This means that what appears as ``drift'' at time $t$ is the accumulated result of all prior inputs, gravity, and interactions.
A system moving fast at time $t$ has high drift accelerations not because of current input, but because of past input that built up the current velocity.


\subsection{Linearization and the Drift--Control Connection}\label{sec:linearization}

Linearization provides a natural bridge from the geometric drift--control decomposition to standard control and estimation workflows.
For control-affine dynamics~\eqref{eq:affine-master}, linearizing at the ZTCF operating point $(x(t), u = 0)$ gives:
\begin{equation}\label{eq:linearized}
\dot{x} \approx \dot{x}_{\mathrm{ZTCF},t} + A_t\,(x - x(t)) + B_t\,u,
\end{equation}
with
\begin{equation}
A_t = \left.\frac{\partial f_{\mathrm{drift}}}{\partial x}\right|_{x(t)}, \qquad B_t = G\bigl(x(t)\bigr).
\end{equation}

This reveals a clear hierarchy:
\begin{itemize}
\item \textbf{ZTCF/drift provides the zeroth-order anchor}: the baseline vector at the current state.
\item \textbf{Linearization adds first-order sensitivity}: how the dynamics change with small perturbations in state and input.
\end{itemize}

The linearized model is the foundation for local stability analysis, feedback design (LQR, pole placement, linear MPC), observer and estimator design, and gain-scheduled controllers along trajectories.


\subsection{Article Layout}

Each calculational framework is treated with the same structure:
\begin{enumerate}
\item General equations in that formalism, with explicit identification of drift and control components.
\item Discussion of the drift--control form, explaining how the affine decomposition appears and what each component represents physically.
\item The calculational method: how computations are organized.
\item A complete double-pendulum derivation.
\item Procedure checklist for implementation.
\end{enumerate}


%% ============================================================
\section{Notation and Modeling Assumptions}\label{sec:notation}
%% ============================================================

Generalized coordinates are $q \in \mathbb{R}^n$, with velocities $\dot{q}$ and accelerations $\ddot{q}$.
The full state is $x = (q, \dot{q}) \in T\mathbb{R}^n$.

Generalized effort $\tau$ denotes torques for rotational coordinates and forces for translational coordinates.
We distinguish control effort $\tau$, constraint reactions $J_c^T \lambda$, and external loads $\tau_{\mathrm{ext}}$.

$A^{-1}$ denotes matrix inverse and $A^T$ denotes transpose.
For unconstrained rigid-body systems, $M(q)$ is symmetric positive definite.

For constrained systems, holonomic constraints are $\phi(q) = 0$ with Jacobian $J_c(q) = \partial\phi/\partial q$.
The velocity-level constraint is $J_c \dot{q} = 0$, and the acceleration-level constraint is $J_c \ddot{q} + \dot{J}_c \dot{q} = 0$.

All ZTCF comparisons are made at fixed instantaneous state unless explicitly noted.


%% ============================================================
\section{Running Example: Planar Double Pendulum}\label{sec:double-pendulum}
%% ============================================================

We use a two-link planar pendulum throughout.
Let
\begin{equation}\label{eq:dp-coords}
q = \begin{bmatrix}\theta_1 \\ \theta_2\end{bmatrix}, \qquad
\dot{q} = \begin{bmatrix}\dot\theta_1 \\ \dot\theta_2\end{bmatrix},
\end{equation}
with link lengths $(l_1,l_2)$, masses $(m_1,m_2)$, and point masses at each distal end.

Endpoint kinematics:
\begin{equation}\label{eq:dp-kinematics}
x_2 = l_1\sin\theta_1 + l_2\sin(\theta_1+\theta_2), \qquad
y_2 = -l_1\cos\theta_1 - l_2\cos(\theta_1+\theta_2).
\end{equation}

The endpoint Jacobian:
\begin{equation}\label{eq:dp-jacobian}
J(q) = \begin{bmatrix}
l_1\cos\theta_1 + l_2\cos(\theta_1+\theta_2) & l_2\cos(\theta_1+\theta_2) \\
l_1\sin\theta_1 + l_2\sin(\theta_1+\theta_2) & l_2\sin(\theta_1+\theta_2)
\end{bmatrix}.
\end{equation}

Manipulator-form dynamics: $\tau = M(q)\ddot{q} + C(q,\dot{q})\dot{q} + g(q)$, with
\begin{equation}\label{eq:dp-mass}
M(q) = \begin{bmatrix}
(m_1+m_2)l_1^2 + m_2l_2^2 + 2m_2l_1l_2\cos\theta_2 & m_2l_2^2+m_2l_1l_2\cos\theta_2 \\
m_2l_2^2+m_2l_1l_2\cos\theta_2 & m_2l_2^2
\end{bmatrix},
\end{equation}
\begin{equation}\label{eq:dp-coriolis}
C(q,\dot{q})\dot{q} = \begin{bmatrix}
-2c\dot\theta_1\dot\theta_2 - c\dot\theta_2^2 \\
c\dot\theta_1^2
\end{bmatrix}, \qquad
c \coloneqq m_2l_1l_2\sin\theta_2,
\end{equation}
\begin{equation}\label{eq:dp-gravity}
g(q) = \begin{bmatrix}
(m_1+m_2)gl_1\sin\theta_1 + m_2gl_2\sin(\theta_1+\theta_2) \\
m_2gl_2\sin(\theta_1+\theta_2)
\end{bmatrix}.
\end{equation}


%% ============================================================
\section{Euler--Lagrange Analysis}\label{sec:euler-lagrange}
%% ============================================================

\subsection{General Equations}

From $L(q,\dot{q}) = T(q,\dot{q}) - V(q)$:
\begin{equation}\label{eq:el-general}
\frac{d}{dt}\left(\frac{\partial L}{\partial \dot{q}}\right) - \frac{\partial L}{\partial q} = \tau
\end{equation}
gives the manipulator equation:
\begin{equation}\label{eq:el-manipulator}
M(q)\,\ddot{q} + C(q,\dot{q})\,\dot{q} + g(q) = \tau,
\end{equation}
where $M(q)\ddot{q}$ is the inertia term, $C(q,\dot{q})\dot{q}$ is the velocity-product (Coriolis/centripetal) term, and $g(q) = \partial V/\partial q$ is the gravity term.

\subsection{The Drift--Control Decomposition}

Solving for acceleration:
\begin{equation}\label{eq:el-drift-control}
\ddot{q} = \underbrace{-M^{-1}(q)\bigl[C(q,\dot{q})\dot{q} + g(q)\bigr]}_{\ddot{q}_{\mathrm{drift}}} + \underbrace{M^{-1}(q)\,\tau}_{\ddot{q}_{\mathrm{ctrl}}}.
\end{equation}

The drift component $\ddot{q}_{\mathrm{drift}} = -M^{-1}(C\dot{q} + g)$ represents the acceleration with zero actuation.
The $-M^{-1}C\dot{q}$ part captures inertial coupling effects; the $-M^{-1}g$ part captures gravitational loading.

The control component $\ddot{q}_{\mathrm{ctrl}} = M^{-1}\tau$ represents the incremental acceleration produced by actuation.
The mass matrix $M^{-1}$ acts as the input gain: it maps applied generalized forces to resulting accelerations.

The drift and control components are structurally independent.
The ZTCF attribution identity holds exactly:
\begin{equation}
\ddot{q}_{\mathrm{total}} - \ddot{q}_{\mathrm{drift}} = M^{-1}\tau = \ddot{q}_{\mathrm{ctrl}}.
\end{equation}

\subsection{Calculational Method}

The Coriolis matrix is obtained from the Christoffel symbols of $M$:
\begin{equation}
C_{ij} = \sum_{k=1}^{n} \frac{1}{2}\left(\frac{\partial M_{ij}}{\partial q_k} + \frac{\partial M_{ik}}{\partial q_j} - \frac{\partial M_{jk}}{\partial q_i}\right)\dot{q}_k.
\end{equation}
At measured state, define $h(q,\dot{q}) \coloneqq C(q,\dot{q})\dot{q} + g(q)$ and compute $\ddot{q}_{\mathrm{drift}} = -M^{-1}h$.

\subsection{Double-Pendulum Derivation}

\subsubsection{Kinetic Energy}

The velocity of mass~1 is $v_1^2 = l_1^2\dot\theta_1^2$.
The velocity of mass~2:
\begin{align}
v_2^2 &= l_1^2\dot\theta_1^2 + l_2^2(\dot\theta_1+\dot\theta_2)^2 + 2l_1l_2\cos\theta_2\,\dot\theta_1(\dot\theta_1+\dot\theta_2).
\end{align}
The total kinetic energy $T = \frac{1}{2}m_1v_1^2 + \frac{1}{2}m_2v_2^2$ yields the mass matrix entries in~\eqref{eq:dp-mass}.

\subsubsection{Potential Energy and Gravity}

$V = -m_1gl_1\cos\theta_1 - m_2g[l_1\cos\theta_1 + l_2\cos(\theta_1+\theta_2)]$.
Taking partial derivatives gives~\eqref{eq:dp-gravity}.

\subsubsection{Coriolis/Centripetal Terms}

From the Christoffel symbols with $c = m_2l_1l_2\sin\theta_2$:
the term $-2c\dot\theta_1\dot\theta_2$ is a Coriolis effect,
$-c\dot\theta_2^2$ is a centripetal effect on joint~1,
and $c\dot\theta_1^2$ is the reaction on joint~2.

\subsubsection{Complete Drift--Control Split}

The drift acceleration is:
\begin{equation}
\ddot{q}_{\mathrm{drift}} = -M^{-1}\begin{bmatrix}
-2c\dot\theta_1\dot\theta_2 - c\dot\theta_2^2 + g_1 \\
c\dot\theta_1^2 + g_2
\end{bmatrix},
\end{equation}
and the control acceleration is $\ddot{q}_{\mathrm{ctrl}} = M^{-1}\tau$.
The total acceleration is the sum: $\ddot{q} = \ddot{q}_{\mathrm{drift}} + \ddot{q}_{\mathrm{ctrl}}$.

For the $2\times 2$ case, $M^{-1}$ has explicit form $M^{-1} = (\det M)^{-1}\bigl[\begin{smallmatrix}M_{22}&-M_{12}\\-M_{21}&M_{11}\end{smallmatrix}\bigr]$.

\subsection{Procedure Checklist}

\begin{enumerate}
\item Define generalized coordinates and physical parameters.
\item Derive $T$ and $V$.
\item Extract $M(q)$, $C(q,\dot{q})$, and $g(q)$.
\item Evaluate drift: $\ddot{q}_{\mathrm{drift}} = -M^{-1}(C\dot{q}+g)$ at measured state.
\item Apply control and compute total acceleration.
\item Compare measured acceleration with model prediction.
\item Use residuals to diagnose unmodeled effects.
\end{enumerate}


%% ============================================================
\section{Constrained Euler--Lagrange Analysis}\label{sec:constrained-el}
%% ============================================================

\subsection{General Equations}

For holonomic constraints $\phi(q)=0$:
\begin{equation}\label{eq:cel-dynamics}
M\ddot{q} + C\dot{q} + g = \tau + J_c^T\lambda + \tau_{\mathrm{ext}},
\end{equation}
\begin{equation}\label{eq:cel-constraint}
J_c\ddot{q} + \dot{J}_c\dot{q} = 0.
\end{equation}

The constraint equation defines an affine hyperplane of admissible accelerations:
\begin{equation}\label{eq:hyperplane}
\mathcal{H}_t \coloneqq \left\{\ddot{q}\;\middle|\;J_c\ddot{q} = -\dot{J}_c\dot{q}\right\}.
\end{equation}

\subsection{The Drift--Control Decomposition with Constraints}

Solving for $\lambda$:
\begin{equation}\label{eq:lambda-general}
\lambda = -(J_cM^{-1}J_c^T)^{-1}\bigl[J_cM^{-1}(\tau+\tau_{\mathrm{ext}}-C\dot{q}-g)+\dot{J}_c\dot{q}\bigr].
\end{equation}

Setting $\tau=0$ gives the constrained drift:
\begin{equation}\label{eq:constrained-drift}
\ddot{q}_{\mathrm{drift}}^{c} = M^{-1}\bigl[-C\dot{q}-g+J_c^T\lambda_0+\tau_{\mathrm{ext}}\bigr].
\end{equation}

The constrained control includes the actuation-induced change in constraint reactions:
\begin{equation}\label{eq:constrained-ctrl}
\ddot{q}_{\mathrm{ctrl}}^{c} = M^{-1}\bigl[\tau+J_c^T(\lambda-\lambda_0)\bigr].
\end{equation}

The decomposition remains additive: $\ddot{q} = \ddot{q}_{\mathrm{drift}}^{c} + \ddot{q}_{\mathrm{ctrl}}^{c}$.

\subsection{The Tangent-Hyperplane Interpretation}

ZTCF gives the base point on $\mathcal{H}_t$.
Control moves along allowable directions:
\begin{equation}\label{eq:hyperplane-control}
\ddot{q} = \ddot{q}_{\mathrm{ZTCF},t}^{c} + B_c(x(t))\,\tau,
\end{equation}
where $B_c = M^{-1}[I - J_c^T(J_cM^{-1}J_c^T)^{-1}J_cM^{-1}]$ is the constraint-consistent input map.
This is exact, not heuristic.

\subsection{Why Constraint Forces Matter Despite Zero Generalized Work}

Even when $P_{\mathrm{constr}}=0$ globally, constraint forces:
\begin{itemize}
\item redirect motion and directly affect which accelerations occur,
\item mediate energy transfer between coordinates (per-coordinate power can be large and nonzero),
\item produce enormous local interface forces (critical in whip-like motions),
\item modify the drift baseline.
\end{itemize}

\subsection{Double-Pendulum Instance: Endpoint on a Circle}

Constrain the endpoint to a circle:
\begin{equation}
\phi(q) = (x_2-x_c)^2 + (y_2-y_c)^2 - R^2 = 0.
\end{equation}
Then $J_c = \partial\phi/\partial q$ and the coupled KKT system
\begin{equation}
\begin{bmatrix}M & -J_c^T \\ J_c & 0\end{bmatrix}
\begin{bmatrix}\ddot{q}\\\lambda\end{bmatrix} =
\begin{bmatrix}\tau-C\dot{q}-g+\tau_{\mathrm{ext}}\\-\dot{J}_c\dot{q}\end{bmatrix}
\end{equation}
yields both constrained accelerations and constraint multipliers.

\subsection{Procedure Checklist}

\begin{enumerate}
\item Define $\phi(q)=0$.
\item Compute $J_c$ and $\dot{J}_c\dot{q}$.
\item Build $\mathcal{H}_t$.
\item Solve the KKT system for $(\ddot{q},\lambda)$.
\item Compute $\ddot{q}_{\mathrm{ZTCF},t}^{c}$ by setting $\tau=0$.
\item Recover $J_c^T\lambda$ and decompose into passive and actuation-induced parts.
\item Compare constrained and unconstrained drift baselines.
\item Audit per-coordinate power flow.
\end{enumerate}


%% ============================================================
\section{Newton--Euler Analysis}\label{sec:newton-euler}
%% ============================================================

\subsection{General Equations}

For each rigid body $i$:
\begin{equation}\label{eq:ne-trans}
m_i a_{c,i} = f_i^{\mathrm{in}} - f_i^{\mathrm{out}} + f_i^{\mathrm{ext}} + f_i^{\mathrm{ctrl}},
\end{equation}
\begin{equation}\label{eq:ne-rot}
I_i\dot\omega_i + \omega_i\times(I_i\omega_i) = n_i^{\mathrm{in}} - n_i^{\mathrm{out}} + n_i^{\mathrm{ext}} + n_i^{\mathrm{ctrl}}.
\end{equation}

Spatial (wrench) form:
\begin{equation}\label{eq:ne-assembled}
\mathcal{M}(x)\dot{v} + \mathcal{C}(x,v)v + \mathcal{G}(x) = \mathcal{W}_{\mathrm{ctrl}} + \mathcal{W}_{\mathrm{constr}} + \mathcal{W}_{\mathrm{ext}}.
\end{equation}

\subsection{The Drift--Control Decomposition}

The body-level drift is the acceleration each body would experience under zero control.
The joint forces under drift ($f_i^{\mathrm{in,drift}}$, $f_i^{\mathrm{out,drift}}$) are constraint reactions---the Newton--Euler analogues of constrained drift forces.

When projected to generalized coordinates through the system Jacobian, the Newton--Euler decomposition reproduces the Euler--Lagrange decomposition exactly.

\subsection{The Recursive Algorithm}

\textbf{Forward recursion (base to tip):} propagate $\omega_i, \dot\omega_i, a_{c,i}$.

\textbf{Backward recursion (tip to base):} compute forces and moments at each joint.
For body $i$:
\begin{equation}
f_i = m_i a_{c,i} - f_i^{\mathrm{ext}} + f_{i+1}, \qquad \tau_i = n_i^T \hat{z}_i.
\end{equation}

ZTCF: set $\tau=0$ and solve the forward-backward system for drift accelerations and passive joint reactions.

\subsection{Double-Pendulum Derivation}

\textbf{Forward recursion:}
Body~1 has $\omega_1=\dot\theta_1$, $\dot\omega_1=\ddot\theta_1$, and COM acceleration $a_{c,1}$.
Body~2 has $\omega_2=\dot\theta_1+\dot\theta_2$, $\dot\omega_2=\ddot\theta_1+\ddot\theta_2$, and COM acceleration $a_{c,2}$ built from the joint~2 acceleration plus rotational terms.

\textbf{Backward recursion:}
Free-body diagram of body~2 yields the joint~2 reaction $f_2^{\mathrm{joint}}$.
Free-body diagram of body~1 uses $-f_2^{\mathrm{joint}}$ (Newton's third law) to determine the joint~1 reaction.

The Newton--Euler formulation makes internal load transfer between bodies explicit: the joint reaction $f_2^{\mathrm{joint}}$ is a specific Cartesian force vector that the Euler--Lagrange form hides inside the mass matrix coupling.

\subsection{Procedure Checklist}

\begin{enumerate}
\item Define local frames, inertias, and COM offsets.
\item Forward kinematics/acceleration recursion.
\item Apply external/contact loads.
\item Backward force-moment recursion.
\item Extract actuator torques and joint reactions.
\item Repeat with $\tau=0$ for ZTCF.
\item Compare drift and total joint reactions.
\item Map local forces to generalized coordinates for cross-validation.
\end{enumerate}


%% ============================================================
\section{Product-of-Exponentials and Screw-Axis Analysis}\label{sec:poe-screw}
%% ============================================================

\subsection{General Equations}

Forward kinematics on $SE(3)$:
\begin{equation}\label{eq:poe-fk}
g_{sb}(q) = e^{\hat{S}_1 q_1}e^{\hat{S}_2 q_2}\cdots e^{\hat{S}_n q_n}g_{sb}(0),
\end{equation}
where $\hat{S}_i\in\mathfrak{se}(3)$ is the twist-hat representation of screw axis $S_i$.

Velocity and acceleration:
\begin{equation}\label{eq:poe-vel}
\mathcal{V}_s = J_s(q)\dot{q}, \qquad
\dot{\mathcal{V}}_s = J_s(q)\ddot{q} + \dot{J}_s(q,\dot{q})\dot{q}.
\end{equation}
The bias acceleration $\dot{J}_s\dot{q}$ is the screw-coordinate analogue of Coriolis/centripetal effects.

Spatial dynamics for each body:
\begin{equation}\label{eq:poe-wrench}
\mathcal{F}_i = \mathcal{I}_i\dot{\mathcal{V}}_i + \mathrm{ad}_{\mathcal{V}_i}^{*}(\mathcal{I}_i\mathcal{V}_i) - \mathcal{F}_i^{\mathrm{ext}}.
\end{equation}

\subsection{The Drift--Control Decomposition}

ZTCF is evaluated as before: freeze state and set actuation to zero.
\begin{equation}
\ddot{q}_{\mathrm{ZTCF},t} = [\ddot{q}(\tau=0)]_{(q(t),\dot{q}(t))}, \qquad
\dot{\mathcal{V}}_{s,\mathrm{ZTCF}} = J_s\ddot{q}_{\mathrm{ZTCF}} + \dot{J}_s\dot{q}.
\end{equation}

The velocity-dependent term $\mathrm{ad}_{\mathcal{V}}^{*}(\mathcal{I}\mathcal{V})$ is the Lie bracket (coadjoint) of the spatial velocity acting on the spatial momentum---the same physics as Coriolis/centripetal terms in Euler--Lagrange, expressed on the Lie algebra.

\subsection{Double-Pendulum Instance}

Screw axes:
\begin{equation}\label{eq:dp-screws}
S_1 = \begin{bmatrix}0\\0\\1\\0\\0\\0\end{bmatrix}, \qquad
S_2 = \begin{bmatrix}0\\0\\1\\0\\-l_1\\0\end{bmatrix}.
\end{equation}

Forward kinematics: $g_{sb}(q)=e^{\hat{S}_1\theta_1}e^{\hat{S}_2\theta_2}g_{sb}(0)$.

Space Jacobian: $J_s = [S_1 \;\; \mathrm{Ad}_{e^{\hat{S}_1\theta_1}}S_2]$, where
\begin{equation}
S_2' = \mathrm{Ad}_{e^{\hat{S}_1\theta_1}}S_2 = \begin{bmatrix}0\\0\\1\\l_1\sin\theta_1\\-l_1\cos\theta_1\\0\end{bmatrix}.
\end{equation}

Extracting the planar components recovers the same endpoint velocity as the geometric Jacobian~\eqref{eq:dp-jacobian}.

\subsection{Procedure Checklist}

\begin{enumerate}
\item Define $g_{sb}(0)$ and screw axes $S_i$.
\item Build PoE forward kinematics and Jacobians.
\item Compute twist and bias acceleration terms.
\item Run recursive Newton--Euler in spatial coordinates.
\item Project wrenches to joint torques.
\item Evaluate ZTCF at fixed state.
\item Compare with Euler--Lagrange ZTCF mapped through the Jacobian.
\end{enumerate}


%% ============================================================
\section{Operational-Space Analysis}\label{sec:operational-space}
%% ============================================================

\subsection{General Equations}

Operational-space analysis reformulates dynamics in terms of \textbf{task variables} that directly describe quantities of interest---endpoint position, orientation, or any function of joint coordinates.

Define task coordinates $x=\psi(q)$ with task Jacobian $J=\partial\psi/\partial q$:
\begin{equation}\label{eq:os-kinematics}
\dot{x}=J\dot{q}, \qquad \ddot{x}=J\ddot{q}+\dot{J}\dot{q}.
\end{equation}

Task-space dynamics:
\begin{equation}\label{eq:os-dynamics}
\Lambda(x)\ddot{x}+\mu(x,\dot{x})+p(x) = F_{\mathrm{ctrl}}+F_{\mathrm{constr}}+F_{\mathrm{ext}},
\end{equation}
where:
\begin{itemize}
\item $\Lambda=(JM^{-1}J^T)^{-1}$ is the \textbf{task-space mass matrix} (operational-space inertia): the effective inertia ``felt'' at the task point.
\item $\mu=\Lambda(JM^{-1}C\dot{q}-\dot{J}\dot{q})$ is the \textbf{task-space Coriolis/centripetal term}.
\item $p=\Lambda JM^{-1}g$ is the \textbf{task-space gravity term}.
\end{itemize}

The mapping between task-space forces and joint torques follows virtual work: $\tau = J^T F$.
For the inverse mapping, the dynamically consistent formulation uses $F = \Lambda J M^{-1}\tau$.

\subsection{The Drift--Control Decomposition}

The joint-level decomposition maps directly to task space.

Task-space drift:
\begin{equation}\label{eq:os-drift}
\ddot{x}_{\mathrm{drift}} = J\ddot{q}_{\mathrm{drift}} + \dot{J}\dot{q} = J[-M^{-1}(C\dot{q}+g)] + \dot{J}\dot{q}.
\end{equation}

This includes both the joint-level drift mapped through the Jacobian ($J\ddot{q}_{\mathrm{drift}}$) and the \textbf{bias acceleration} $\dot{J}\dot{q}$ due to the changing Jacobian.

Task-space control: $\ddot{x}_{\mathrm{ctrl}} = JM^{-1}\tau$.

Total: $\ddot{x} = \ddot{x}_{\mathrm{drift}} + \ddot{x}_{\mathrm{ctrl}}$.

The drift answers ``how would the task point accelerate with all actuators off?''
If task-level constraints are active, the same tangent-hyperplane interpretation applies in $\ddot{x}$-space.

\subsection{Double-Pendulum Instance}

For endpoint task $x=[x_2\;y_2]^T$, the task Jacobian is~\eqref{eq:dp-jacobian} and the task-space inertia is $\Lambda=(JM^{-1}J^T)^{-1}$, a $2\times 2$ configuration-dependent matrix.

At stretched configurations ($\theta_2\approx 0$), the endpoint inertia along the radial direction is high.
At folded configurations ($\theta_2\approx\pi$), the kinematic mapping changes character.

The task-space drift is:
\begin{equation}
\ddot{x}_{\mathrm{drift}} = J\ddot{q}_{\mathrm{drift}} + \dot{J}\dot{q},
\end{equation}
where $\ddot{q}_{\mathrm{drift}}=-M^{-1}(C\dot{q}+g)$ from the Euler--Lagrange computation.

Singularities occur when $\theta_2=0$ or $\theta_2=\pi$ (collinear links), where $J$ loses rank and $\Lambda$ becomes singular---infinite force would be required to produce acceleration along the lost direction.

\subsection{Procedure Checklist}

\begin{enumerate}
\item Choose task variables $x=\psi(q)$.
\item Derive $J$ and $\dot{J}\dot{q}$.
\item Compute $\Lambda,\mu,p$.
\item Specify desired task-space acceleration or force.
\item Map task commands to joint torques via $J^T$ and the null-space projection.
\item Evaluate drift baseline in task space: $\ddot{x}_{\mathrm{drift}}=J\ddot{q}_{\mathrm{drift}}+\dot{J}\dot{q}$.
\item Compare total and drift task-space accelerations.
\item Check for singularities and examine task-space inertia variation.
\end{enumerate}


%% ============================================================
\section{Coordinate-Free Geometric Analysis}\label{sec:geometric}
%% ============================================================

\subsection{General Equations}

On configuration manifold $Q$ with kinetic-energy Riemannian metric:
\begin{equation}\label{eq:geo-eom}
\nabla_{\dot{q}}\dot{q} + \mathrm{grad}\,V(q) = \mathcal{G}(q)u + \mathcal{F}_{\mathrm{ext}} + \mathcal{F}_{\mathrm{constr}}.
\end{equation}

Control-affine state form on $TQ$:
\begin{equation}\label{eq:geo-affine}
\dot{x} = f_{\mathrm{drift}}(x) + \sum_{i=1}^{m}g_i(x)u_i.
\end{equation}

\subsection{The Drift--Control Decomposition}

Intrinsic drift:
\begin{equation}\label{eq:geo-drift}
f_{\mathrm{drift}}(q,\dot{q}) = \begin{pmatrix}\dot{q}\\-\Gamma(\dot{q},\dot{q})-\mathrm{grad}\,V(q)+\mathcal{F}_{\mathrm{ext}}\end{pmatrix},
\end{equation}
where $\Gamma(\dot{q},\dot{q})$ is the Christoffel connection applied to velocity (the same physics as $M^{-1}C\dot{q}$ in Euler--Lagrange).

Intrinsic control: $g_i(q,\dot{q}) = (0, \mathcal{G}_i(q))$.

The geometric perspective makes transparent that:
the drift field is intrinsic (independent of coordinates),
the connection encodes inertial coupling,
constraints restrict to submanifolds,
and ZTCF is geodesic deviation modified by potential.

\subsection{Double-Pendulum Instance}

Configuration manifold: $Q = S^1\times S^1$ (a torus).

The kinetic-energy metric in coordinates $(\theta_1,\theta_2)$ is the mass matrix~\eqref{eq:dp-mass}.
The Christoffel symbols generate the Coriolis/centripetal terms.
With zero potential and zero control, the equations reduce to the geodesic equation:
\begin{equation}
\ddot\theta_i + \Gamma^i_{jk}\dot\theta_j\dot\theta_k = 0.
\end{equation}

The full drift field on $TQ$ matches the Euler--Lagrange drift exactly.
Coordinate independence is verified by switching to $(\phi_1,\phi_2)=(\theta_1,\theta_1+\theta_2)$: the predicted physical accelerations are identical despite different coordinate expressions.

\subsection{Procedure Checklist}

\begin{enumerate}
\item Define $Q$ and admissible states on $TQ$.
\item Build the kinetic-energy metric.
\item Compute the connection and potential gradient.
\item Add control and external covector fields.
\item Simulate intrinsic flow with and without control.
\item Use coordinate charts only for computation/plotting.
\item Verify conclusions by changing charts.
\end{enumerate}


%% ============================================================
\section{Cross-Formalism Equivalence}\label{sec:cross-formalism}
%% ============================================================

\subsection{The Same Physics, Different Languages}

For the same mechanism at the same state, each formalism encodes the same force balance:
\begin{align}
\text{Euler--Lagrange:}\quad & M\ddot{q}+C\dot{q}+g-J_c^T\lambda = \tau+\tau_{\mathrm{ext}}, \label{eq:xf-gen}\\
\text{Screw:}\quad & \mathcal{I}\dot{\mathcal{V}}+\mathrm{ad}_{\mathcal{V}}^*(\mathcal{I}\mathcal{V})+\mathcal{W}_g-\mathcal{W}_{\mathrm{constr}} = \mathcal{W}_{\mathrm{ctrl}}+\mathcal{W}_{\mathrm{ext}}, \label{eq:xf-screw}\\
\text{Operational-space:}\quad & \Lambda\ddot{x}+\mu+p = F_{\mathrm{ctrl}}+F_{\mathrm{constr}}+F_{\mathrm{ext}}, \label{eq:xf-os}\\
\text{Geometric:}\quad & \nabla_{\dot{q}}\dot{q}+\mathrm{grad}\,V = \mathcal{G}u+\mathcal{F}_{\mathrm{ext}}+\mathcal{F}_{\mathrm{constr}}. \label{eq:xf-geo}
\end{align}

\subsection{Term-by-Term Correspondence}

\textbf{Inertia:} $M\ddot{q}$ $\leftrightarrow$ $\mathcal{I}\dot{\mathcal{V}}$ $\leftrightarrow$ $\Lambda\ddot{x}$ $\leftrightarrow$ inertial part of $\nabla_{\dot{q}}\dot{q}$.

\textbf{Velocity-product:} $C\dot{q}$ $\leftrightarrow$ $\mathrm{ad}_{\mathcal{V}}^*(\mathcal{I}\mathcal{V})+\dot{J}_s\dot{q}$ $\leftrightarrow$ $\mu$ $\leftrightarrow$ $\Gamma(\dot{q},\dot{q})$.

\textbf{Gravity:} $g(q)$ $\leftrightarrow$ gravity wrenches $\leftrightarrow$ $p(x)$ $\leftrightarrow$ $\mathrm{grad}\,V$.

\textbf{Constraints:} $J_c^T\lambda$ $\leftrightarrow$ constraint wrenches $\leftrightarrow$ $F_{\mathrm{constr}}$ $\leftrightarrow$ reaction covectors.

\textbf{Control:} $\tau$ $\leftrightarrow$ control wrenches $\leftrightarrow$ $F_{\mathrm{ctrl}}$ $\leftrightarrow$ $\mathcal{G}u$.

\subsection{Frame Changes and Invariants}

Twists and wrenches transform as $V_s=\mathrm{Ad}_g V_b$, $F_s=\mathrm{Ad}_g^{-T}F_b$.
Instantaneous power is frame-invariant: $F_s^TV_s = F_b^TV_b$.

\subsection{ZTCF Anchor in Each Formalism}

At fixed state:
\begin{itemize}
\item Euler--Lagrange: evaluate $\ddot{q}$ with $\tau=0$.
\item Newton--Euler: run recursion with control removed.
\item Screw: evaluate spatial dynamics with zero control wrenches.
\item Operational-space: map passive acceleration to task space.
\item Geometric: evaluate intrinsic drift field.
\end{itemize}

Constraints do not break equivalence: all frameworks predict the same physical accelerations, forces, and power flows at any state.


%% ============================================================
\section{Hidden Loads and Internal Energy Transfer}\label{sec:hidden-loads}
%% ============================================================

Compact equations hide joint reactions, internal energy transfer, actuator internals, and distributed contact loads.

\subsection{Power-Channel Decomposition}

\begin{equation}
P_{\mathrm{ctrl}}=\tau^T\dot{q}, \qquad
P_{\mathrm{drift}}=(C\dot{q}+g)^T\dot{q}, \qquad
P_{\mathrm{constr}}=(J_c^T\lambda)^T\dot{q}=0.
\end{equation}

Even when total constraint power is zero, per-coordinate power $P_{\mathrm{constr},i}=[J_c^T\lambda]_i\dot{q}_i$ can be large---only their sum is zero.
Constraint forces actively extract power from some coordinates and inject it into others.

\subsection{Recommended Workflow}

\begin{enumerate}
\item Add explicit constraint and external-load channels.
\item Solve the augmented KKT system for $(\ddot{q},\lambda)$.
\item Map generalized reactions to local force/torque quantities.
\item Audit global and per-coordinate/per-body power flow.
\end{enumerate}


%% ============================================================
\section{Symmetry, Invariance, and Conservation Laws}\label{sec:symmetry}
%% ============================================================

Symmetry underlies physical law.
The equations of motion are objective: they do not depend on the coordinate frame.
Coordinate changes alter symbolic form, not physical content.

By Noether's theorem:
\begin{itemize}
\item Translation symmetry $\Rightarrow$ linear momentum conservation.
\item Rotation symmetry $\Rightarrow$ angular momentum conservation.
\item Time-translation symmetry $\Rightarrow$ energy conservation.
\end{itemize}

Conservation laws interact with the drift--control decomposition:
if a quantity is conserved under drift, then control is the only mechanism that can change it.
If energy is conserved under drift, the drift dynamics preserve energy surfaces and control moves between them.

These are structural properties of the physics, not of one coordinate system or framework.
They hold in every representation and provide framework-independent cross-validation checks.


%% ============================================================
\section{Workflow Comparison}\label{sec:workflow}
%% ============================================================

\begin{center}
\small
\renewcommand{\arraystretch}{1.2}
\begin{tabularx}{\textwidth}{>{\raggedright\arraybackslash}p{3.1cm}Y Y Y}
\toprule
Formalism & Best question answered & Core unknowns & Practical workflow \\
\midrule
Euler--Lagrange & How inputs change generalized accelerations & $\ddot{q}$ or $\tau$ & Derive $M,C,g$; compute drift; compare with controlled response \\
\midrule
Constrained E--L & How constraints redirect motion and loads & $(\ddot{q},\lambda)$ & Solve KKT system; recover $J_c^T\lambda$ \\
\midrule
Newton--Euler & Where local loads and transfer pathways occur & Body forces, moments, joint torques & Forward/backward recursion; interpret load paths \\
\midrule
PoE / screw & How motion, twists, wrenches compose across frames & $g_{sb}$, $J$, twists, wrenches & PoE kinematics; recursive propagation \\
\midrule
Operational-space & How task objectives map to joint effort & Task accel/force, joint torques & Build $J,\Lambda,\mu,p$; close task loop \\
\midrule
Coordinate-free & Which conclusions are representation-invariant & Intrinsic drift and input fields & Formulate on $Q$/$TQ$; verify across charts \\
\bottomrule
\end{tabularx}
\end{center}


%% ============================================================
\section{Conclusion}\label{sec:conclusion}
%% ============================================================

This article has demonstrated three interlocking principles through six calculational frameworks:

\textbf{The drift--control decomposition is universal.}
Every framework admits a clean affine split $\dot{x}=f_{\mathrm{drift}}(x)+G(x)u$.
In Euler--Lagrange form: $\ddot{q}=-M^{-1}(C\dot{q}+g)+M^{-1}\tau$.
In Newton--Euler: per-body force balance with and without actuator loads.
In screw form: wrench decomposition.
In operational-space: task-level drift and control.
In geometric form: intrinsic drift on $TQ$.

\textbf{Constraint forces are active participants.}
Despite zero net generalized work, constraint forces redirect motion, mediate energy transfer, and modify the drift baseline.
The tangent-hyperplane interpretation provides exact geometric meaning.

\textbf{Representation independence is the physics.}
The equations are objective.
Different frameworks produce equivalent predictions.
Conservation laws are framework-independent.
Cross-validation between methods is possible and powerful.

The ZTCF provides a unifying diagnostic across all frameworks: at any instant, it answers ``what would the system do without actuation?''
This single question isolates state effects from control, defines the constraint hyperplane anchor, and connects to linearization for design.

\end{document}
